Desenhistas utilizam linhas de fuga para orientar seus traços no papel. Linhas de fuga são linhas traçadas de tal modo a constituir um espaço tridimensional no plano bidimensional da folha em que se desenha, e raramente fazem parte do desenho acabado, sendo apagadas pelo desenhista quando não precisa mais delas. Penso em Cecinepá como uma linha de fuga pois se trata de um mundo traçado no espaço-tempo presente -- daí seu caráter heterotópico, e não utópico -- a fim de se oferecerem parâmetros que orientem a composição de um novo mundo comum, no qual Cecinepá se mostre obsoleto.